\documentclass[a4paper]{article}
\usepackage[utf8]{inputenc}
\usepackage{graphicx}
\usepackage{amsmath}
\usepackage[bottom=2.0cm,top=2.0cm,left=2.0cm,right=2.0cm]{geometry}
\usepackage[portuges]{babel}
\usepackage{indentfirst}
\usepackage{hyperref}  %%%%
\hypersetup{colorlinks,citecolor=black,filecolor=black,linkcolor=black,urlcolor=black} %%%%


\begin{document}

%%%%% Com o aumento do número de professores que pedem relatórios em formato digital, com imagens coloridas e etc, me vi motivado a montar este template para jogar na núvem e compartilhar a formatação que mais usei/uso nesses meus anos de UERJ.

%%%%% Não é de minha responsabilidade se algum professor não aceitar essa formatação, fiz este modelo SEM qualquer compromisso com normas de formatação da universidade, o fiz apenas para facilitar a confecção de relatórios. Alguns professores exigem uma formatação específica, preste atenção.

\begin{titlepage}
	\begin{center}
		\begin{figure}[htb!]
			\begin{flushleft}
				\includegraphics[width=3.9cm,height=3.9cm]{imagens/uerj.jpg}
			\end{flushleft}
		\end{figure}
        \vspace{-3.8cm}
        \hspace{2.1cm}\Large{\textbf{Universidade do Estado do Rio de Janeiro}}\\
        \hspace{2.1cm}\Large{Faculdade de Engenharia}\\
        \hspace{2.1cm}\Large{Departamento da disciplina}\\
        \hspace{2.1cm}\Large{Disciplina}\\
        
        \vspace{200pt}
        
        \LARGE{\textbf{Relatório xxxXXXxxx:}}\\
        \Large{experimento}\\
        
        \vspace{200pt}
        
        \hfill \underline{Aluno:}\\
        \hfill \large{xxxxXXXXXxxxx}\\
        \vspace{25pt}
        \hfill \underline{\Large{Professor:}}\\
        \hfill \large{xxxxxxXXXXXXxxxxxx}\\
        
        
        \vspace{\fill}
        mês\\
        ano
        
	\end{center}
\end{titlepage}


%% o índice remissivo pode ser retirado se excluir o texto desde esse comentário até a próxima sequência de %, não se esqueça também de excluir os hyperref lá no alto, estarão marcados com simbolos.

\newpage

\Large\tableofcontents
\thispagestyle{empty}
             %%%%%
\newpage
\pagenumbering{arabic}
\large
\section{Introdução}
\subsection{Teoria}
\subsection{Objetivo}

\section{Preparatório}
\subsection{Cálculos}
\subsection{Simulação}

\section{Experimento}

\section{Análise de Dados}

\section{Conclusão}

\section{Referências}

\end{document}
